\section{Proof sketches (OC/ETS-constrained)}
\label{sec:proofs}

All arguments in this section are proof sketches in the strict sense.
They derive logical consequences directly from ETS predicates, the fixed
precedence ordering, and the phase-level allowance relation.
No empirical data, dynamical assumptions, control semantics, or temporal
models are invoked.

\subsection{Uniqueness of phase classification}

\begin{proof}[Sketch of Proposition~\ref{prop:unique-phase}]
Let $s \in \Omega(K_{EA})$ be arbitrary.
By construction, each ETS phase predicate
$\pi_{\texttt{PHASE}}(s)$ is Boolean-valued.
The phase logic defines a total precedence ordering
\[
\texttt{STOP} \succ \texttt{COLLAPSE} \succ \texttt{INERTIA} \succ \texttt{SUCCESS}.
\]
If exactly one predicate is true at $s$, the corresponding phase is assigned.
If multiple predicates are simultaneously true, the precedence ordering
selects a unique maximal element.
Hence the classification map $\Pi(s)$ is well-defined and single-valued
for all admissible states.
\end{proof}

\subsection{Discontinuity at threshold crossings}

\begin{proof}[Sketch of Proposition~\ref{prop:discontinuous}]
Each phase predicate is defined by strict Boolean conditions on ETS variables,
including inequalities of the form $f_k(s) > 0$ and Boolean flags such as
$\texttt{observability\_lost}(s)$.
These predicates induce a partition of $\Omega(K_{EA})$ into disjoint
subsets separated by threshold hypersurfaces.
A change in predicate truth value corresponds to crossing such a boundary
and induces a jump between distinct phase classes.
Since ETS imposes no continuity or smoothness requirements on state
parametrizations, the induced classification map $\Pi(s)$ may change
discontinuously.
\end{proof}

\subsection{Absorbing nature of \texttt{STOP}}

\begin{proof}[Sketch of Proposition~\ref{prop:stop-absorbing}]
The ETS allowance relation specifies no outgoing phase-level transitions
from \texttt{STOP}. Formally,
\[
\texttt{STOP} \Rightarrow \varnothing .
\]
Therefore, once a state is classified as \texttt{STOP}, no further
phase-level transitions are admissible under ETS.
\end{proof}

\subsection{Irreversibility of \texttt{COLLAPSE}}

\begin{proof}[Sketch of Proposition~\ref{prop:collapse-no-return}]
According to the ETS allowance relation, the only admissible outgoing
transition from \texttt{COLLAPSE} is
\[
\texttt{COLLAPSE} \Rightarrow \{\texttt{STOP}\}.
\]
No transitions from \texttt{COLLAPSE} to either \texttt{SUCCESS} or
\texttt{INERTIA} are permitted.
Hence return to non-terminal phases is excluded at the level of phase logic.
\end{proof}

\subsection{Non-steerability of phase transitions}

\begin{proof}[Sketch of Proposition~\ref{prop:non-steerable}]
ETS specifies phase transitions as a guarded allowance relation rather than
as a dynamical system with trajectories, controls, or admissible paths.
Consequently, the existence of phase boundaries or allowed transitions
does not imply reachability, controllability, or reversibility of state
evolution within $\Omega(K_{EA})$.
\end{proof}

\subsection{Scope of the proof sketches}

\begin{remark}
The proof sketches above establish only logical consequences of the ETS
structure. They make no claims about temporal evolution, causal mechanisms,
recovery feasibility, or empirical realizability, all of which lie outside
the scope of this paper.
\end{remark}
